%%=============================================================================
%% Methodologie
%%=============================================================================

\chapter{\IfLanguageName{dutch}{Methodologie}{Methodology}}%
\label{ch:methodologie}

%% In dit hoofstuk geef je een korte toelichting over hoe je te werk bent
%% gegaan. Verdeel je onderzoek in grote fasen, en licht in elke fase toe wat
%% de doelstelling was, welke deliverables daar uit gekomen zijn, en welke
%% onderzoeksmethoden je daarbij toegepast hebt. Verantwoord waarom je
%% op deze manier te werk gegaan bent.
%% 
%% Voorbeelden van zulke fasen zijn: literatuurstudie, opstellen van een
%% requirements-analyse, opstellen long-list (bij vergelijkende studie),
%% selectie van geschikte tools (bij vergelijkende studie, "short-list"),
%% opzetten testopstelling/PoC, uitvoeren testen en verzamelen
%% van resultaten, analyse van resultaten, ...
%%
%% !!!!! LET OP !!!!!
%%
%% Het is uitdrukkelijk NIET de bedoeling dat je het grootste deel van de corpus
%% van je bachelorproef in dit hoofstuk verwerkt! Dit hoofdstuk is eerder een
%% kort overzicht van je plan van aanpak.
%%
%% Maak voor elke fase (behalve het literatuuronderzoek) een NIEUW HOOFDSTUK aan
%% en geef het een gepaste titel.

Het onderzoek wordt uitgevoerd in vier fasen, die samen een gestructureerd en iteratief plan van aanpak vormen. Elke fase draagt bij aan het ontwikkelen en evalueren van een proof-of-concept voor security monitoring en incident response binnen een SaaS- en DevOps-context. De fasen bouwen bewust op elkaar voort; de probleemanalyse stuurt het architectuurontwerp, het ontwerp vormt de basis voor de implementatie, en de evaluatie toetst het geheel aan de vooropgestelde doelstellingen en deelvragen.

\section{Fase 1: Probleemanalyse}
In de eerste fase wordt in samenwerking met Devinity de huidige situatie grondig onderzocht. De bestaande monitoringtools, werkwijze en knelpunten worden in kaart gebracht via observaties, gesprekken met collega's en een eerste verkenning van beschikbare logs en systemen. Ook worden bestaande incidenten en de manier waarop deze momenteel worden opgevolgd geanalyseerd. Die analyse is niet enkel beschrijvend, het doel is een scherp en onderbouwd beeld te krijgen van waar de grootste pijnpunten zich bevinden, zodat het architectuurontwerp in de volgende fase daar gericht op kan inspelen.

Meer bepaald wordt onderzocht:

\begin{itemize}
    \item welke uitdagingen Devinity ervaart bij het gebruik van meerdere tools voor security monitoring (deelvraag P1);
    
    \item hoe incidenten momenteel worden gedetecteerd en gerapporteerd, en welke tekortkomingen daarin bestaan op vlak van snelheid, correlatie en consistentie (deelvraag P2);
    
    \item welke securitydata beschikbaar is vanuit de bestaande systemen en hoe deze momenteel worden gecentraliseerd (deelvraag P3);
    
    \item welke metrics en KPI's vandaag worden gebruikt en waar deze tekortschieten (deelvraag P4);
    
    \item hoe de huidige aanpak de operationele efficiëntie en het risicomanagement beïnvloedt (deelvraag P5).
\end{itemize}

Het resultaat van deze fase is een gedetailleerde probleemstelling die de basis vormt voor het architectuurontwerp. Hierbij wordt ook expliciet rekening gehouden met de niet-disruptieve doelstelling. De analyse brengt niet alleen tekortkomingen in kaart, maar ook welke bestaande tools en werkwijzen behouden en verder benut kunnen worden.

\section{Fase 2: Architectuurontwerp}
Op basis van de probleemanalyse wordt een architectuur ontworpen voor de centrale verzameling en verwerking van securitydata. Verschillende technologieën worden geselecteerd op basis van hun aansluiting bij de bestaande infrastructuur van Devinity, hun schaalbaarheid en de mate waarin ze integreerbaar zijn zonder grote aanpassingen aan de bestaande werking. De datastromen worden gedefinieerd via pipelines zoals API-polling, webhooks en logforwarders, en voor elke keuze wordt expliciet verantwoord waarom die boven alternatieven de voorkeur krijgt. De architectuur bouwt voort op de kernprincipes van SIEM-systemen, namelijk logaggregatie, normalisatie, en rapportering, en houdt rekening met de NIST-richtlijnen voor intrusion detection en prevention ~\autocite{Scarfone_Mell_2007}.

Concreet wordt in deze fase bepaald:

\begin{itemize}
    \item welke architectuur en technologieën het meest geschikt zijn voor een geautomatiseerd securitymonitoringsysteem binnen een DevOps-omgeving (deelvraag O1);
    
    \item welke data-integratieprocessen het meest efficiënt zijn om logs en securitydata te verzamelen en te normaliseren (deelvraag O3);
    
    \item hoe de beveiliging en betrouwbaarheid van het systeem zelf kunnen worden gegarandeerd binnen de bestaande pipeline (deelvraag O6);
    
    \item welke KPI's later in de evaluatiefase zullen worden gebruikt (deelvraag O5).
\end{itemize}

Een belangrijke voorwaarde van dit ontwerp is dat de oplossing zo weinig mogelijk ingrijpt op de bestaande infrastructuur. Ontwerpkeuzes worden daarom bewust afgestemd op de tools en werkwijzen die al aanwezig zijn binnen Devinity. Het resultaat van deze fase is een gedocumenteerd systeemontwerp met een beschrijving van de gekozen architectuur, de databronnen, de ingestiepipeline en het centrale analyseplatform.

\section{Fase 3: Proof-of-concept ontwikkeling}
In deze fase wordt het ontwerp omgezet naar een werkend prototype binnen de geselecteerde omgeving. Kernfuncties zoals logverzameling, regelgebaseerde detectie, eenvoudige anomaly-detectie en automatische rapportage worden stap voor stap geïmplementeerd en getest. Die iteratieve aanpak laat toe om tussentijds bij te sturen wanneer bepaalde ontwerpkeuzes in de praktijk anders uitpakken dan verwacht.

Testen worden uitgevoerd met gesimuleerde data en ongevoelige logdata uit testomgevingen, bestaande uit authenticatielogs, auditlogs, API-events en doelbewust gegenereerde security-events zoals mislukte logins, om realistische scenario's te kunnen simuleren zonder productiedata te gebruiken of de live omgeving te verstoren.

Zo wordt uitgewerkt hoe rule-based detectie op afwijkende HTTP-statuscodes, als aanvulling op regelgebaseerde detectie, kan bijdragen aan het identificeren van mogelijk ongeautoriseerde toegangspogingen (deelvraag O2), en hoe automatische rapportgeneratie kan worden geïmplementeerd zodat rapporten zowel technische details als managementinzichten bevatten (deelvraag O4).

Tegelijk wordt in de praktijk getoetst of de gekozen data-integratieprocessen effectief werken zoals voorzien in het ontwerp (deelvraag O3).

Het resultaat van deze fase is een werkend proof-of-concept dat de centralisatie, analyse en rapportage van securitydata aantoont binnen de omgeving van Devinity.

\section{Fase 4: Evaluatie}
In de laatste fase worden de prestaties van het proof-of-concept geëvalueerd aan de hand van vooraf gedefinieerde KPI's, zoals detectiesnelheid, responstijd en de verhouding tussen relevante en irrelevante meldingen. Die KPI's werden niet ad hoc vastgelegd, maar bewust gekozen op basis van de tekortkomingen die in de probleemanalyse naar voren kwamen, zodat de evaluatie een directe vergelijking mogelijk maakt tussen de oude en de nieuwe situatie. De metingen worden uitgevoerd over meerdere testrondes om toevallige uitschieters te minimaliseren en een betrouwbaar beeld te geven van de gemiddelde en maximale prestaties van het systeem.

Deze fase toetst in welke mate de geïmplementeerde oplossing de eerder vastgestelde tekortkomingen effectief aanpakt en oplost. Zo wordt getoetst of incidenten sneller en consistenter worden gedetecteerd (deelvraag P2), of de gecentraliseerde aanpak zorgt voor een beter overzicht van securitydata (deelvraag P3), en of de gedefinieerde KPI's een objectieve meting van de effectiviteit mogelijk maken (deelvraag P4 en O5). Ook wordt beoordeeld of de oplossing de operationele efficiëntie verbetert ten opzichte van de huidige situatie (deelvraag P5).

Het resultaat van deze fase is een evaluatierapport met conclusies en aanbevelingen voor verdere uitbouw van het systeem binnen de infrastructuur van Devinity.

\begin{figure}[H]
    \centering
    \includegraphics[width=\textwidth]{methodologie_diagram.png}
    \caption[Overzicht van de methodologie in vier fasen.]{\label{fig:methodologie}
        Schematisch overzicht van de vierfasen-methodologie van dit onderzoek.}
\end{figure}

De aanpak bouwt voort op de NIST-aanbevelingen voor intrusion detection en prevention systemen~\autocite{Scarfone_Mell_2007}, waarin het vastleggen en analyseren van security-events centraal staat voor effectieve detectie en respons. Tegelijkertijd baseert de architectuurkeuze zich op kerncomponenten van SIEM-systemen, zoals gecentraliseerde logverzameling, normalisatie en event-monitoring~\autocite{s21144759}.